%------------------------------------------------------------------------------%
\documentclass[fleqn,utf8,aspectratio=169,12pt]{beamer}

\usepackage{gaal}

\title{Escalonamento de Matrizes}

%------------------------------------------------------------------------------%
\begin{document}

\addtitle

%------------------------------------------------------------------------------%
\section{Sistemas Equivalentes}

%------------------------------------------------------------------------------%
\begin{frame}{Sistemas Equivalentes}
  Dois sistemas lineares são equivalentes
  \\ \pause
  quando possuem as mesmas soluções

  \bigskip
  \pause
  Operações elementares sobre linhas produzem sistemas equivalentes
\end{frame}

%------------------------------------------------------------------------------%
\begin{frame}{Operações Elementares Sobre Linhas}
\begin{enumerate}[<+->]
    \item Troca de duas linhas
    \[
      L_i\leftrightarrow L_j
    \]

    \item Multiplicação de uma linha por um número não nulo
    \[
      L_i\leftarrow kL_i,
      \qquad k\neq0
    \]

    \item Adição de um múltiplo de uma linha a outra
    \[
      L_i\leftarrow L_i+kL_j
    \]
\end{enumerate}
\end{frame}

%------------------------------------------------------------------------------%
%------------------------------------------------------------------------------%
\begin{question}
Operação elementar: \emph{troca de linhas}

Considerando
\(
  A =
  \begin{bmatrix}
    0 & 2 & 1 \\
    3 & 4 & 5 \\
    1 & 2 & 3
  \end{bmatrix}
\)

Troque a primeira e a segunda linhas
\[
  L_1\leftrightarrow L_2
\]
\end{question}

%------------------------------------------------------------------------------%
\begin{resolution}{Solução}
\[
\begin{bmatrix}
  0 & 2 & 1 \\
  3 & 4 & 5 \\
  1 & 2 & 3
\end{bmatrix}
\]

\pause
Aplicando
\(
  L_1\leftrightarrow L_2
\)

\pause
\[
\begin{bmatrix}
  \emph{3} & \emph{4} & \emph{5} \\
  \emph{0} & \emph{2} & \emph{1} \\
  1 & 2 & 3
\end{bmatrix}
\]
\end{resolution}

%------------------------------------------------------------------------------%

%------------------------------------------------------------------------------%
\begin{question}
Operação elementar: \emph{multiplicação de uma linha}

Considerando
\(
  A =
  \begin{bmatrix}
    2 & 4 & 6 \\
    1 & 3 & 5
  \end{bmatrix}
\)

Transforme o primeiro elemento da primeira linha em 1, fazendo
\[
  L_1 \leftarrow \frac{1}{2} L_1
\]
\end{question}

%------------------------------------------------------------------------------%
\begin{resolution}{Solução}
\[
  \begin{bmatrix}
  2 & 4 & 6 \\
  1 & 3 & 5
  \end{bmatrix}
\]

\pause
Nova linha 1
\[
  L'_1
  \pause
  \;=\;
  \frac{1}{2} L_1
  \pause
  \;=\;
  \frac{1}{2}
  \begin{bmatrix}
    2 & 4 & 6 \\
  \end{bmatrix}
  \pause
  \;=\;
  \begin{bmatrix}
    1 & 2 & 3 \\
  \end{bmatrix}
\]

\pause
\[
  \begin{bmatrix}
  1 & 2 & 3 \\
  1 & 3 & 5
  \end{bmatrix}
\]
\end{resolution}

%------------------------------------------------------------------------------%

%------------------------------------------------------------------------------%
\begin{question}
  Operação elementar: \emph{combinação de linhas}

  Considerando
  \(
    A =
    \begin{pmatrix}
      1 & 2 & 3 \\
      2 & 5 & 7
    \end{pmatrix}
  \)

  Produza um zero abaixo do primeiro pivô, usando
  \[
    L_2\leftarrow L_2-2L_1
  \]
\end{question}

%------------------------------------------------------------------------------%
\begin{resolution}{Solução}
\[
  \begin{pmatrix}
    1 & 2 & 3 \\
    2 & 5 & 7
  \end{pmatrix}
\]

\pause
\[
  L'_2
  \pause
  \;=\;
   L_2 - 2L_1
   \pause
  \;=\;
  \begin{pmatrix}
    2 & 5 & 7
  \end{pmatrix}
  -
  2\begin{pmatrix}
    1 & 2 & 3
  \end{pmatrix}
  \pause
  \;=\;
  \begin{pmatrix}
    0 & 1 & 1
  \end{pmatrix}
\]

\pause
\[
  \begin{pmatrix}
    1 & 2 & 3 \\
    0 & 1 & 1
  \end{pmatrix}
\]
\end{resolution}

%------------------------------------------------------------------------------%


%------------------------------------------------------------------------------%
\section{Escalonamento de Matrizes}

%------------------------------------------------------------------------------%
\begin{frame}{Por que escalonar uma matriz?}

\onslide<+->{Transformar uma matriz em uma forma \emph{mais simples}} \\
\onslide<+->{preservando informações importantes}

\medskip

\onslide<+->{Aplicações}
\begin{itemize}[<+->]
    \item resolução de sistemas lineares
    \item determinação do posto de uma matriz
    \item identificação da existência e quantidade de soluções
    \item cálculo da inversa de uma matriz
    \item análise de dependência linear
\end{itemize}

\medskip
\onslide<+->{%
Utilizar operações elementares para produzir
\emph{zeros em posições estratégicas}
}
\end{frame}

%------------------------------------------------------------------------------%
\begin{frame}{Matriz Escalonada}
\onslide<+->{Uma matriz está em \textbf{forma escalonada} quando:}
\begin{enumerate}[<+->]
    \item todas as linhas nulas estão abaixo das linhas não nulas

    \item o \emph{pivo} (primeiro elemento não nulo da linha)
    está à direita do pivo da linha anterior

    \item todos os elementos abaixo de um pivô são nulos
\end{enumerate}

\medskip
\onslide<+->{Exemplo de matriz escalonada
\qquad
\(
\begin{pmatrix}
  \emph{1} & 2        & 1        & 3  \\
  \z       & \emph{1} & 4        & 2  \\
  \z       & \z       & \emph{2} & 5  \\
  \z       & \z       & \z       & \z
\end{pmatrix}
\)}
\end{frame}

%------------------------------------------------------------------------------%
%------------------------------------------------------------------------------%
\begin{question}
Determine quais das matrizes abaixo estão em forma escalonada
\[
  A =
  \begin{pmatrix}
    1 & 2 & 3 \\
    0 & 1 & 4 \\
    0 & 0 & 2
  \end{pmatrix}
  \qquad
  B =
  \begin{pmatrix}
    1 & 2 & 3 \\
    0 & 0 & 4 \\
    0 & 2 & 5
  \end{pmatrix}
  \qquad
  C =
  \begin{pmatrix}
    0 & 1 & 3 \\
    0 & 0 & 2 \\
    0 & 0 & 0
  \end{pmatrix}
\]
\end{question}

%------------------------------------------------------------------------------%
\begin{resolutiont}{Solução}
\onslide<+->{\[
  A = \begin{pmatrix}
    1 & 2 & 3 \\
    0 & 1 & 4 \\
    0 & 0 & 2
  \end{pmatrix}
\]}
\onslide<+->{A matriz $A$ está escalonada}
\begin{enumerate}[<+->]
  \item os pivôs estão progressivamente mais à direita
  \item todos os elementos abaixo dos pivôs são nulos
\end{enumerate}
\end{resolutiont}

%------------------------------------------------------------------------------%
\begin{resolutiont}{Solução}
\onslide<+->{\[
  B = \begin{pmatrix}
    1 & 2                   & 3                   \\
    0 & 0                   & {\color<3->{red} 4} \\
    0 & {\color<3->{red} 2} & 5
  \end{pmatrix}
\]}%
\onslide<+->{A matriz $B$ não está escalonada}  \\
\onslide<+->{o pivô da terceira linha está à esquerda do pivô da segunda linha}
\end{resolutiont}

%------------------------------------------------------------------------------%
\begin{resolutiont}{Solução}
\onslide<+->{\[
  C = \begin{pmatrix}
    0 & 1 & 3 \\
    0 & 0 & 2 \\
    0 & 0 & 0
  \end{pmatrix}
\]}
\onslide<+->{A matriz $C$ está escalonada}
\begin{enumerate}[<+->]
  \item os pivôs estão progressivamente mais à direita
  \item todos os elementos abaixo dos pivôs são nulos
  \item as linhas nulas estão abaixo das linhas não nulas
\end{enumerate}
\end{resolutiont}

%------------------------------------------------------------------------------%

%------------------------------------------------------------------------------%
\begin{question}
Escalone a matriz
\(
  A=
  \begin{bmatrix}
    1  & 2  & -1 \\
    2  & 5  & 1  \\
    -1 & -1 & 2
  \end{bmatrix}
\)
\end{question}

%------------------------------------------------------------------------------%
\begin{resolution}{Solução}
  Utilizamos operações elementares sobre as linhas\\
  até obter uma matriz em forma escalonada
\end{resolution}

%------------------------------------------------------------------------------%
\begin{resolution}{Primeira Coluna}
\begin{columns}
\column{0.3\linewidth}
  \[
  \begin{bmatrix}
    1  & 2  & -1 \\
    2  & 5  & 1  \\
    -1 & -1 & 2
  \end{bmatrix}
  \]
  \pause
\column{0.7\linewidth}
  Eliminamos os elementos abaixo do primeiro pivô
  \\ \pause
  \(L_2 \leftarrow L_2-2 L_1\)
  \\ \pause
  \(L_3\leftarrow L_3+L_1\)
\end{columns}
\vspace{-\baselineskip}
\[
  \pause
  L'_2
  \;=\;
  L_2 - 2 L_1
  \pause
  \;=\;
  \begin{bmatrix}
    2  & 5  & 1
  \end{bmatrix}
  - 2
  \begin{bmatrix}
    1  & 2  & -1
  \end{bmatrix}
  \pause
  \;=\;
  \begin{bmatrix}
    0 & 1 & 3
  \end{bmatrix}
\]
\[
  \pause
  L'_3
  \;=\;
  L_3 + L_1
  \pause
  \;=\;
  \begin{bmatrix}
    -1 & -1 & 2
  \end{bmatrix}
  +
  \begin{bmatrix}
    1  & 2  & -1
  \end{bmatrix}
  \pause
  \;=\;
  \begin{bmatrix}
    0 & 1 & 1
  \end{bmatrix}
\]
\pause
\[
  \begin{bmatrix}
    1 & 2 & -1 \\
    0 & 1 & 3  \\
    0 & 1 & 1
  \end{bmatrix}
\]
\end{resolution}

%------------------------------------------------------------------------------%
\begin{resolution}{Segunda Coluna}
\begin{columns}
\column{0.3\linewidth}
  \[
    \begin{bmatrix}
      1 & 2 & -1 \\
      0 & 1 & 3  \\
      0 & 1 & 1
    \end{bmatrix}
  \]
  \pause
\column{0.7\linewidth}
  Eliminamos os elementos abaixo do segundo pivô
  \\ \pause
  \(L_3\leftarrow L_3-L_2\)
\end{columns}
\vspace{-\baselineskip}
\[
  \pause
  L'_3
  \;=\;
  L_3 - L_2
  \pause
  \;=\;
  \begin{bmatrix}
      0 & 1 & 1
  \end{bmatrix}
  -
  \begin{bmatrix}
    0 & 1 & 3
  \end{bmatrix}
  \pause
  \;=\;
  \begin{bmatrix}
    0 & 0 & -2
  \end{bmatrix}
\]

\pause
\[
  \begin{bmatrix}
    1 & 2 & -1 \\
    0 & 1 & 3  \\
    0 & 0 & -2
  \end{bmatrix}
\]
\end{resolution}

%------------------------------------------------------------------------------%
\begin{resolution}{Solução}
Portanto, uma forma escalonada de $A$ é
\(
  \begin{bmatrix}
    1 & 2 & -1 \\
    0 & 1 & 3  \\
    0 & 0 & -2
  \end{bmatrix}
\)
\end{resolution}

%------------------------------------------------------------------------------%
\begin{resolution}{Notação Resumida}
\[
\begin{NiceArray}{rrrl}[columns-width=3ex]
    1 &  2 & -1  & \\
    2 &  5 &  1  & \\
   -1 & -1 &  2  &
  \CodeAfter
    \SubMatrix[{1-1}{3-3}]
  \end{NiceArray}
\]

\pause
\[
  \begin{NiceArray}{rrrl}[columns-width=3ex]
    \hpm 1 & 2 & -1 & \\
         0 & 1 &  3 & \quad L_2 \leftarrow L_2-2 L_1 \\
         0 & 1 &  1 & \quad L_3 \leftarrow L_3+L_1
  \CodeAfter
    \SubMatrix[{1-1}{3-3}]
  \end{NiceArray}
\]

\pause
\[
  \begin{NiceArray}{rrrl}[columns-width=3ex]
    \hpm 1 & 2 & -1 & \\
         0 & 1 &  3 & \\
         0 & 0 & -2 & \quad L_3\leftarrow L_3-L_2
  \CodeAfter
    \SubMatrix[{1-1}{3-3}]
  \end{NiceArray}
\]
\end{resolution}

%------------------------------------------------------------------------------%

%------------------------------------------------------------------------------%
\begin{question}
Escalone a matriz
\(
  A =
  \begin{pmatrix}
    1  & 2  & 1 \\
    2  & 5  & 4 \\
    -1 & -1 & 1
  \end{pmatrix}
\)
\end{question}

%------------------------------------------------------------------------------%
\begin{resolution}{Solução}

Partimos de
\[
\begin{pmatrix}
1&2&1\\
2&5&4\\
-1&-1&1
\end{pmatrix}.
\]
Aplicamos
\[
L_2\leftarrow L_2-2L_1
\]
e
\[
L_3\leftarrow L_3+L_1.
\]
Obtemos
\[
\begin{pmatrix}
1&2&1\\
0&1&2\\
0&1&2
\end{pmatrix}.
\]
Agora fazemos
\[
L_3\leftarrow L_3-L_2.
\]
Portanto,
\[
\begin{pmatrix}
1&2&1\\
0&1&2\\
0&0&0
\end{pmatrix}.
\]
\end{resolution}



%------------------------------------------------------------------------------%
\section{Escalonamento de Sistemas Lineares}

%------------------------------------------------------------------------------%
\begin{frame}{Sistemas Lineares}
O sistema
\[
  \begin{cases}
    a_{11} x + a_{12} y + a_{13} z = b_1 \\
    a_{21} x + a_{22} y + a_{23} z = b_2 \\
    a_{21} x + a_{23} y + a_{33} z = b_3
  \end{cases}
\]
\pause
Pode ser representado matricialmente como
\[
  A {\bm x} = {\bm b}
\]
\pause
onde
\[
  A =
  \begin{bmatrix}
    a_{11} & a_{12} & a_{13} \\
    a_{21} & a_{22} & a_{23} \\
    a_{21} & a_{23} & a_{33}
  \end{bmatrix}
  \pause
  \qquad
  {\bm x} =
  \begin{bmatrix}
    x \\ y \\ z
  \end{bmatrix}
  \pause
  \qquad
  {\bm b} =
  \begin{bmatrix}
    b_1 \\ b_2 \\ b_3
  \end{bmatrix}
\]
\end{frame}

%------------------------------------------------------------------------------%
\begin{frame}{Matriz Aumentada}
  \emph{Matriz Aumentada} do sistema \(A{\bm x} = {\bm b}\)
  \pause
  \[
    \; [ \, A \, | \, {\bm b} \, ]
  \]
  \[
    \pause
    \left[\begin{array}{rrr|r}
      a_{11} & a_{12} & a_{13} & b_1 \\
      a_{21} & a_{22} & a_{23} & b_2 \\
      a_{21} & a_{23} & a_{33} & b_3
    \end{array}\right]
  \]

  \pause
  \medskip
  Ao escalonar a matriz aumentada obtemos \\
  uma versão simplificada do sistema linear original
\end{frame}

%------------------------------------------------------------------------------%
%------------------------------------------------------------------------------%
\begin{question}
  Utilize o escalonamento para resolver o sistema linear
  \[
    \systeme{
       x + 2y -  z = 3,
      2x + 5y +  z = 8,
      -x - y  + 2z = 1
      }
  \]
\end{question}

%------------------------------------------------------------------------------%
\begin{resolution}{Matriz Aumentada}
\[
  A = \left[
  \begin{array}{rrr}
    1  & 2  & -1  \\
    2  & 5  & 1   \\
    -1 & -1 & 2
  \end{array}
  \right]
  \pause
  \qquad\qquad
  b = \left[
  \begin{array}{rrr|r}
     3 \\
     8 \\
     1
  \end{array}
  \right]
\]

\pause
Matriz aumentada
\[
  \left[
  \begin{array}{rrr|r}
    1  & 2  & -1 & 3 \\
    2  & 5  & 1  & 8 \\
    -1 & -1 & 2  & 1
  \end{array}
  \right]
\]
\end{resolution}

%------------------------------------------------------------------------------%
\begin{resolution}{Escalonamento do Sistema}
\begin{columns}
\column{0.4\linewidth}
  \[\;\;\,
    \left[
    \begin{array}{rrr|r}
      1  & 2  & -1 & 3 \\
      2  & 5  & 1  & 8 \\
      -1 & -1 & 2  & 1
    \end{array}
    \right]
  \]
  \pause
\column{0.6\linewidth}
  Eliminar os elementos abaixo do primeiro pivô
  \\ \pause
  \(
    L_2 \leftarrow L_2 - 2 L_1
  \)
  \\ \pause
  \(
    L_3 \leftarrow L_3 + L_1
  \)
\end{columns}
\vspace{-\baselineskip}
\[
  \pause
  L'_2
  \;=\;
  L_2 - 2 L_1
  \pause
  \;=\;
  \left[
  \begin{array}{rrr|r}
    2  & 5  & 1  & 8
  \end{array}
  \right]
  -2
  \left[
  \begin{array}{rrr|r}
    1  & 2  & -1 & 3
  \end{array}
  \right]
  \pause
  \;=\;
  \left[
  \begin{array}{rrr|r}
    0 & 1 & 3  & 2
  \end{array}
  \right]
\]
\[
  \pause
  L'_3
  \;=\;
  L_3 + L_1
  \pause
  \;=\;
  \left[
  \begin{array}{rrr|r}
    -1 & -1 & 2  & 1
  \end{array}
  \right]
  +
  \left[
  \begin{array}{rrr|r}
    1  & 2  & -1 & 3
  \end{array}
  \right]
  \pause
  \;=\;
  \left[
  \begin{array}{rrr|r}
    0 & 1 & 1  & 4
  \end{array}
  \right]
\]

\pause
\[
  \left[
  \begin{array}{rrr|r}
    1 & 2 & -1 & 3 \\
    0 & 1 & 3  & 2 \\
    0 & 1 & 1  & 4
  \end{array}
  \right]
\]
\end{resolution}

%------------------------------------------------------------------------------%
\begin{resolution}{Escalonamento do Sistema}
\begin{columns}
\column{0.4\linewidth}
  \[\;\;\,
    \left[
    \begin{array}{rrr|r}
      1  & 2  & -1 & 3 \\
      2  & 5  & 1  & 8 \\
      -1 & -1 & 2  & 1
    \end{array}
    \right]
  \]
  \pause
\column{0.6\linewidth}
  Eliminar os elementos abaixo do segundo pivô
  \\ \pause
  \(
    L_3 \leftarrow L_3 - L_2
  \)
\end{columns}
\vspace{-\baselineskip}
\[
  \pause
  L'_3
  \;=\;
  L_3 - L_2
  \pause
  \;=\;
  \left[
  \begin{array}{rrr|r}
    -1 & -1 & 2  & 1
  \end{array}
  \right]
  -
  \left[
  \begin{array}{rrr|r}
    2  & 5  & 1  & 8
  \end{array}
  \right]
  \pause
  \;=\;
  \left[
  \begin{array}{rrr|r}
    0 & 0 & -2 & 2
  \end{array}
  \right]
\]

\pause
\[
  \left[
  \begin{array}{rrr|r}
    1 & 2 & -1 & 3 \\
    0 & 1 & 3  & 2 \\
    0 & 0 & -2 & 2
  \end{array}
  \right]
\]
\end{resolution}

%------------------------------------------------------------------------------%
\begin{resolution}{Sistema Linear Escalonado}
Matriz aumentada do sistema escalonado
\[
  \left[
  \begin{array}{rrr|r}
    1  & 2  & -1 & 3 \\
    \z & 1  & 3  & 2 \\
    \z & \z & -2 & 2
  \end{array}
  \right]
\]

\pause
Sistema linear escalonado
\[
  \systeme{
    x + 2y  -z = 3,
         y +3z = 2,
           -2z = 2}
\]
\end{resolution}

%------------------------------------------------------------------------------%
\begin{resolution}{Substituição Retroativa}
\begin{columns}[t]
\column{0.26\linewidth}
\onslide<+->{\centerline{Última linha}}
\begin{align*}
  \onslide<+->{-2z & = 2}  \\
  \onslide<+->{z   & = -1}
\end{align*}
\column{0.27\linewidth}
\onslide<+->{\centerline{Segunda linha}}
\begin{align*}
  \onslide<+->{y + 3z    & = 2} \\
  \onslide<+->{y + 3(-1) & = 2} \\
  \onslide<+->{y - 3     & = 2} \\
  \onslide<+->{y         & = 5}
\end{align*}
\column{0.4\linewidth}
\onslide<+->{\centerline{Primeira linha}}
\begin{align*}
  \onslide<+->{x + 2y - z          & = 3}  \\
  \onslide<+->{x + 2\times 5 -(-1) & = 3}  \\
  \onslide<+->{x + 10 + 1          & = 3}  \\
  \onslide<+->{x                   & = -8}
\end{align*}
\end{columns}
\end{resolution}

%------------------------------------------------------------------------------%
\begin{resolution}{Solução}
A solução do sistema é
\(
  {\bm x}
  \;=\;
  \begin{bmatrix}
    x \\ y \\ z
  \end{bmatrix}
  \;=\;
  \begin{bmatrix}
    -8 \\ 5 \\ -1
  \end{bmatrix}
\)
\end{resolution}

%------------------------------------------------------------------------------%

%------------------------------------------------------------------------------%
\begin{question}
  Resolva o sistema, utilizando escalonamento,
  % \[
  %   \systeme{
  %      x +  y + z = 6,
  %     2x -  y + z = 3,
  %      x + 2y - z = 2}
  % \]
  \[
    \left\{\begin{array}{rrrr}
      x  & +  y & + z & = 6 \\
      2x & -  y & + z & = 3 \\
      x  & + 2y & - z & = 2
    \end{array}\right.
  \]
\end{question}

%------------------------------------------------------------------------------%
\begin{resolution}{Matriz Aumentada}
  Matriz aumentada
  \[
    \begin{amatrix}{3}
      1 & 1  & 1  & 6 \\
      2 & -1 & 1  & 3 \\
      1 & 2  & -1 & 2
    \end{amatrix}
  \]
\end{resolution}

%------------------------------------------------------------------------------%
\begin{resolution}{Escalonamento do Sistema}
\begin{columns}
\column{0.4\linewidth}
  \[\;\;\,
    \begin{amatrix}{3}
      1 & 1  & 1  & 6 \\
      2 & -1 & 1  & 3 \\
      1 & 2  & -1 & 2
    \end{amatrix}
  \]
  \pause
\column{0.6\linewidth}
  Eliminar os elementos abaixo do primeiro pivô
  \\ \pause
  \(
    L_2 \leftarrow L_2 - 2 L_1
  \)
  \\ \pause
  \(
    L_3 \leftarrow L_3 - L_1
  \)
\end{columns}
\vspace{-\baselineskip}
\[
  \pause
  L'_2
  \;=\;
  L_2 - 2 L_1
  \pause
  \;=\;
  \begin{amatrix}{3}
    2 & -1 & 1  & 3
  \end{amatrix}
  - 2
  \begin{amatrix}{3}
    1 & 1  & 1  & 6
  \end{amatrix}
  \pause
  \;=\;
  \begin{amatrix}{3}
    0 & -3 & -1 & -9
  \end{amatrix}
\]
\[
  \pause
  L'_3
  \;=\;
    L_3 - L_1
  \pause
  \;=\;
  \begin{amatrix}{3}
    1 & 2  & -1 & 2
  \end{amatrix}
  -
  \begin{amatrix}{3}
    1 & 1  & 1  & 6
  \end{amatrix}
  \pause
  \;=\;
  \begin{amatrix}{3}
    0 & 1  & -2 & -4
  \end{amatrix}
\]
\pause
\[
  \begin{amatrix}{3}
    1 & 1  & 1  & 6  \\
    0 & -3 & -1 & -9 \\
    0 & 1  & -2 & -4
  \end{amatrix}
\]
\end{resolution}

%------------------------------------------------------------------------------%
\begin{resolution}{Escalonamento do Sistema}
\begin{columns}
\column{0.4\linewidth}
  \[\;\;\,
  \begin{amatrix}{3}
      1 & 1  & 1  & 6  \\
      0 & -3 & -1 & -9 \\
      0 & 1  & -2 & -4
  \end{amatrix}
  \]
  \pause
\column{0.6\linewidth}
  Trocamos a segunda e terceira linhas
  \\
  \(
    L_2\leftrightarrow L_3
  \)
\end{columns}
\vspace{-\baselineskip}
\[
  \pause
  L'_2
  \;=\;
  L_3
  \pause
  \;=\;
  \begin{amatrix}{3}
    0 & 1  & -2 & -4
  \end{amatrix}
\]
\[
  \pause
  L'_3
  \;=\;
  L_2
  \pause
  \;=\;
  \begin{amatrix}{3}
    0 & -3 & -1 & -9
  \end{amatrix}
\]
\pause
\[
  \begin{amatrix}{3}
    1 & 1  & 1  & 6  \\
    0 & 1  & -2 & -4 \\
    0 & -3 & -1 & -9
  \end{amatrix}
\]
\end{resolution}

%------------------------------------------------------------------------------%
\begin{resolution}{Escalonamento do Sistema}
\begin{columns}
\column{0.4\linewidth}
  \[\;\;\,
  \begin{amatrix}{3}
      1 & 1  & 1  & 6  \\
      0 & 1  & -2 & -4 \\
      0 & -3 & -1 & -9
  \end{amatrix}
  \]
  \pause
\column{0.6\linewidth}
  Eliminar os elementos abaixo do segundo pivô
  \\ \pause
  \(
    L_3 \leftarrow L_3 + 3 L_2
  \)
\end{columns}
\vspace{-\baselineskip}
\[
  \pause
  L'_3
  \;=\;
    L_3 + 3 L_2
  \pause
  \;=\;
  \begin{amatrix}{3}
    0 & -3 & -1 & -9
  \end{amatrix}
  + 3
  \begin{amatrix}{3}
    0 & 1  & -2 & -4
  \end{amatrix}
\]
\[
  \pause
  \hphantom{L'_3}
  \;=\;
  \begin{amatrix}{3}
    0 & 0 & -7 & -21
  \end{amatrix}
\]
\pause
\[
  \begin{amatrix}{3}
    1 & 1 & 1  & 6   \\
    0 & 1 & -2 & -4  \\
    0 & 0 & -7 & -21
  \end{amatrix}
\]
\end{resolution}

%------------------------------------------------------------------------------%
\begin{resolution}{Sistema Linear Escalonado}
  Matriz aumentada do sistema escalonado
  \[
    \begin{amatrix}{3}
      1  & 1  & 1  & 6   \\
      \z & 1  & -2 & -4  \\
      \z & \z & -7 & -21
    \end{amatrix}
  \]
  
  \pause
  Sistema linear escalonado
  % \[
  %   \systeme{
  %     x + y +  z =  6,
  %         y - 2z = -4,
  %           - 7z = -21}
  % \]
  \[
    \left\{\begin{array}{rrrr}
      x & + y & +  z & =  6 \\
        &   y & - 2z & = -4 \\
        &     & - 7z & = -21
    \end{array}\right.
  \]
\end{resolution}

%------------------------------------------------------------------------------%
\begin{resolution}{Substituição Retroativa}
 \begin{columns}[t]
 \column{0.26\linewidth}
 \onslide<+->{\centerline{Última linha}}
 \begin{align*}
   \onslide<+->{-7 z & = -21} \\
   \onslide<+->{z    & = 3}
 \end{align*}
 \column{0.27\linewidth}
 \onslide<+->{\centerline{Segunda linha}}
 \begin{align*}
   \onslide<+->{y -2 z        & = -4} \\
   \onslide<+->{y -2 \times 3 & = -4} \\
   \onslide<+->{y -6          & = -4} \\
   \onslide<+->{y             & = 2}
 \end{align*}
 \column{0.4\linewidth}
 \onslide<+->{\centerline{Primeira linha}}
 \begin{align*}
   \onslide<+->{x + y + z & = 6} \\
   \onslide<+->{x + 2 + 3 & = 6} \\
   \onslide<+->{x         & = 1}
 \end{align*}
 \end{columns}
\end{resolution}

%------------------------------------------------------------------------------%
\begin{resolution}{Solução}
A solução do sistema é
\(
  \mathbf{x}
  \;=\;
  \begin{bmatrix}
    x \\ y \\ z
  \end{bmatrix}
  \;=\;
  \begin{bmatrix}
    1 \\ 2 \\ 3
  \end{bmatrix}
\)
\end{resolution}

%------------------------------------------------------------------------------%


%------------------------------------------------------------------------------%
%\section{Lista Mínima}
%
%%------------------------------------------------------------------------------%
%\begin{frame}{Lista Mínima}
%  \begin{enumerate}
%    \item
%  \end{enumerate}
%
%  \vfill
%  Atenção: A prova é baseada no livro, não nas apresentações
%\end{frame}

%------------------------------------------------------------------------------%
\end{document}
%------------------------------------------------------------------------------%
